\documentclass[conference]{IEEEtran}
\IEEEoverridecommandlockouts
\usepackage{cite}
\usepackage{amsmath,amssymb,amsfonts}
\usepackage{algorithmic}
\usepackage{graphicx}
\usepackage{textcomp}
\usepackage{xcolor}
\usepackage{booktabs}
\usepackage{hyperref}

% Define the image placeholders environment
\newcommand{\imageplaceholder}[2]{
    \begin{figure}[h]
    \centering
    \includegraphics[width=0.42\textwidth]{#1}
    \caption{#2}
    \label{fig:#1}
    \end{figure}
}

\def\BibTeX{{\rm B\kern-.05em{\sc i\kern-.025em b}\kern-.08em
    T\kern-.1667em\lower.7ex\hbox{E}\kern-.125emX}}

\begin{document}

\title{Visual Odometry for SLAM Preprocessing in Monocular Camera Systems}

\author{
\IEEEauthorblockN{\textit{Guided by:} Prof. Mohandas R}
\IEEEauthorblockA{
Associate Professor \\
\textit{Cambridge Institute of Technology} \\
mohandas.eee@cambridge.edu.in
}
\and
\IEEEauthorblockN{Roshani Bankar}
\IEEEauthorblockA{
\textit{Cambridge Institute of Technology} \\
roshani.22eee@cambridge.edu.in
}
}


\maketitle

\begin{abstract}
Autonomous navigation in disaster response environments demands fast, reliable and drift-resistant localization systems capable of operating in GPS-denied, visually degraded, and structurally hazardous environments. This paper presents a deeply expanded study of a hybrid localization architecture combining \textbf{Monocular Visual Odometry (VO)} and \textbf{Extended Kalman Filter-based Simultaneous Localization and Mapping (EKF-SLAM)}. VO provides high-frequency feature-based motion estimates derived from Lucas–Kanade optical flow and Essential Matrix decomposition. Crucially, the EKF integrates these high-rate VO predictions with low-rate absolute corrections from wheel odometry, laser scan matching, and artificial \textbf{RF landmarks} for global drift correction and scale recovery. This extended paper includes a comprehensive literature review, detailed algorithmic derivations, operational analysis, and a rigorous discussion of disaster-environment constraints including dust, occlusion, and debris. The results demonstrate that VO improves the prediction accuracy of EKF, while landmark updates correct the scale and trajectory drift, achieving robust performance suitable for real-world search-and-rescue robotics.
\end{abstract}

\begin{IEEEkeywords}
SLAM, Visual Odometry, EKF, Disaster Robotics, Sensor Fusion, Autonomous Navigation, Scale Correction
\end{IEEEkeywords}

% ====================================================
\section{Introduction}

Disaster response robotics has rapidly evolved to support rescue and reconnaissance operations in highly \textbf{hazardous and unstructured environments}, such as collapsed buildings, underground tunnels, and fire-damaged industrial facilities. In such scenarios, conventional \textbf{Global Positioning System (GPS) is entirely unavailable} or unreliable. Furthermore, the operational environment is characterized by \textbf{poor and dynamic lighting conditions}, \textbf{dust/smoke obscuring sensor readings}, and a topology that undergoes \textbf{continuous change}. The core challenge for autonomous robots in these settings is robust \textbf{Simultaneous Localization and Mapping (SLAM)}.

Classical SLAM approaches relying on LiDAR or wheel odometry are highly vulnerable under disaster conditions: \textbf{LiDAR} suffers from significant backscatter in smoke, and \textbf{wheel odometry} accumulates excessive, non-Gaussian drift on rubble and debris.

\textbf{Monocular cameras} offer a compelling, lightweight, and low-cost alternative. \textbf{Visual Odometry (VO)}, which estimates the camera's motion by tracking visual features between successive frames, has become crucial for modern SLAM systems. However, VO alone is plagued by two fundamental issues: \textbf{scale ambiguity} (motion is recovered only up to a scale factor) and \textbf{cumulative drift}.

To address these limitations, this work proposes and implements a \textbf{hybrid localization architecture} that tightly integrates high-frequency Monocular VO as a robust pose \textit{prediction} source within an \textbf{Extended Kalman Filter-based SLAM (EKF-SLAM)} pipeline. The EKF provides a powerful, probabilistic framework to fuse the VO estimates with auxiliary sensor data—specifically, laser scan matching, wheel odometry, and absolute measurements from artificial \textbf{RF beacons (landmarks)}—for global drift and scale correction.

This expanded paper significantly enlarges the original work by integrating additional methodologies, comparisons, technical analysis, and an extended set of scholarly references.

% ---
\section{Literature Survey}

The field of SLAM is rooted in probabilistic robotics, with \textbf{Thrun et al. \cite{Thrun2005}} providing the foundational reasoning. Early foundational algorithms included the \textbf{Extended Kalman Filter (EKF-SLAM)} and the Particle Filter-based FastSLAM. Later research transitioned toward graph-based optimization and \textbf{Bundle Adjustment (BA)} \cite{Triggs2000}.

\textbf{Juneja’s} early work on EKF-SLAM in MATLAB \cite{Juneja2013} and textured-scene localization \cite{Juneja2014} highlighted the necessity of robust landmark management and visual features. \textbf{Scaramuzza and Fraundorfer \cite{scaramuzza2011}} provided a comprehensive review of Visual Odometry (VO) techniques, stressing the trade-off between robustness and computational cost. \textbf{Mur-Artal’s ORB-SLAM \cite{mur2015}} demonstrated keyframe-based methods, introducing robust loop closures.

In the specialized domain of \textbf{Disaster Robotics}, studies by pioneers like \textbf{Murphy \cite{murphy2014}} and \textbf{Tadokoro \cite{tadokoro2009}}, and contemporary work by \textbf{Queralta et al. \cite{queralta2020}}, consistently stress the paramount need for \textbf{multi-sensor redundancy} against unique environmental hazards. Direct VO pipelines such as LSD-SLAM \cite{engel2014} and DSO \cite{engel2017} highlight failure modes under low texture, emphasizing the need for EKF fusion.

Recent advancements in hybrid systems, such as \textbf{VINS-Mono \cite{qin2018}}, have inspired the multi-modal design. Other foundational SLAM analyses \cite{forster2014, cadena2016} consistently conclude that robust localization for field robotics requires strategic sensor fusion to overcome individual sensor weaknesses. This body of research strongly supports the hybrid VO–EKF design proposed in this paper.

% ---
\section{Objectives}

The objectives of this expanded work are defined to address the specific localization challenges encountered by search-and-rescue robots in GPS-denied, visually degraded disaster environments:

\begin{itemize}
    \item To analytically investigate and implement \textbf{Monocular VO} as the primary, high-rate motion \textbf{preprocessing step} for the EKF-SLAM pose prediction phase.
    \item To implement and parameterize Lucas–Kanade optical flow and \textbf{Essential Matrix decomposition} (using RANSAC) for robust relative camera motion estimation.
    \item To design a \textbf{robust sensor fusion strategy} within the EKF-SLAM framework that selectively fuses the VO-derived pose estimates with auxiliary sensor data: wheel odometry, laser scan matching (ICP), and absolute position updates from artificial RF beacons (landmarks).
    \item To demonstrate that the EKF's measurement updates, particularly from RF beacons, can correct the \textbf{cumulative scale and trajectory drift} inherent to Monocular VO.
    \item To evaluate the system's stability and accuracy in simulated disaster conditions, explicitly considering scenarios involving \textbf{dust, severe occlusion}, and \textbf{non-Gaussian sensor noise}.
    \item To provide a scalable, modular architecture suitable for real-time deployment in search-and-rescue robots.
\end{itemize}

% ---
\section{Methodology}

The proposed system integrates a high-frequency \textbf{Monocular Visual Odometry (VO) module} with a low-frequency, robust \textbf{Extended Kalman Filter (EKF-SLAM) architecture}.

\subsection{System Architecture}
The hybrid pipeline operates asynchronously:

$$
\text{Camera Feed} \xrightarrow{\text{High Rate}} \text{VO Prediction}
$$
$$
\text{VO Prediction} \xrightarrow{\substack{\text{Low Rate} \\ \text{or Triggered}}} \text{EKF Correction} \rightarrow \text{Pose \& Map}
$$

The VO Module outputs a relative pose change $(\Delta R, \Delta \mathbf{t})$ which serves as the control input ($\mathbf{u}_k$) for the EKF's prediction step. The EKF Correction step is triggered by auxiliary sensor measurements.

\subsection{Monocular Visual Odometry}
The VO module follows a classical feature-based pipeline.

\subsubsection{Feature Extraction}
\textbf{Shi–Tomasi corners} are used for stable tracking. The corner response function $R$ is defined by the eigenvalues $\lambda_i$ of the gradient covariance matrix $\mathbf{G}$:
$$
R = \min(\lambda_1, \lambda_2)
$$
where
$$
\mathbf{G} = \sum_{x,y} \mathbf{w}(x,y) \begin{bmatrix} I_x^2 & I_x I_y \\ I_x I_y & I_y^2 \end{bmatrix}
$$

\subsubsection{Lucas–Kanade Optical Flow}
Features are tracked between frame $I_k$ and $I_{k+1}$ using the iterative \textbf{Lucas–Kanade pyramidical optical flow} method, minimizing the sum-of-squared differences (SSD) based on the \textbf{brightness constancy assumption}:
$$
I_k(x,y) \approx I_{k+1}(x+u, y+v)
$$
The displacement $\mathbf{d} = [u, v]^T$ is computed iteratively. \textbf{Pyramidal implementation} is used to handle large inter-frame motion.

\subsubsection{Essential Matrix Estimation}
The 2D-to-2D feature correspondences are used to compute the \textbf{Essential Matrix ($E$)}, which encodes the relative motion $(\Delta R, \Delta \mathbf{t})$. The \textbf{epipolar constraint} must be satisfied:
$$
\mathbf{x}_{k+1}^T E \mathbf{x}_k = 0
$$
The matrix $E$ is estimated using the \textbf{8-point algorithm} combined with \textbf{RANSAC} to robustly eliminate outliers. $E$ is then decomposed via \textbf{SVD} to recover the relative rotation $\Delta R$ and translation direction $\Delta \mathbf{t}$.

\subsubsection{Trajectory Integration and Scale}
The relative motion is integrated over time:
$$
T_k = T_{k-1} \cdot \begin{bmatrix} \Delta R & \Delta \mathbf{t} \\ \mathbf{0}^T & 1 \end{bmatrix}
$$
The unknown absolute scale $s_k$ is factored in, $\mathbf{t}_k = s_k \cdot \Delta \mathbf{t}$. This scale factor is dynamically estimated and corrected externally by the EKF updates, using known-distance measurements (e.g., between two RF beacons).

% --- New Section V with algorithmic detail
\section{Algorithmic Derivations for EKF Fusion}

The core of the hybrid system lies in the probabilistic management of uncertainty via the EKF. This section details the prediction step using VO input and the correction step using landmark observations.

\subsection{State Prediction (VO Input)}
The state vector $\mathbf{x}_{k-1}$ is transformed into the predicted state $\mathbf{\hat{x}}_k$ using the nonlinear motion model $\mathbf{f}$, driven by the scaled VO output $\mathbf{u}_{vo} = [s_k \Delta x, s_k \Delta y, \Delta \theta]^T$, where $s_k$ is the current best scale estimate. For a 2D robot pose $[x, y, \theta]^T$:
\begin{align*}
    \mathbf{\hat{x}}_{r,k} &= \mathbf{f}(\mathbf{x}_{r,k-1}, \mathbf{u}_{vo}) \\
    &= \begin{bmatrix}
        x_{k-1} + s_k \Delta x \cos(\theta_{k-1}) - s_k \Delta y \sin(\theta_{k-1}) \\
        y_{k-1} + s_k \Delta x \sin(\theta_{k-1}) + s_k \Delta y \cos(\theta_{k-1}) \\
        \theta_{k-1} + \Delta \theta
    \end{bmatrix}
\end{align*}
The covariance $\mathbf{P}_{k-1}$ is propagated to $\mathbf{\hat{P}}_k$ using the Jacobian $\mathbf{F}_k$ of the motion model $\mathbf{f}$ with respect to the state, and the noise covariance matrix $\mathbf{Q}_k$:
$$
\mathbf{\hat{P}}_k = \mathbf{F}_{k-1} \mathbf{P}_{k-1} \mathbf{F}_{k-1}^T + \mathbf{Q}_k
$$
The matrix $\mathbf{Q}_k$ is crucial as its large elements relating to $x$ and $y$ reflect the high positional uncertainty and unbounded scale error inherent in monocular VO.

\subsection{Measurement Correction (RF Beacons)}
The EKF correction phase is triggered by an observation $\mathbf{z}_k$ of an RF beacon (landmark $i$) at global coordinates $\mathbf{m}_i = [m_x, m_y]^T$. The observed measurement is $\mathbf{z}_k = [r, \alpha]^T$, where $r$ is the measured range and $\alpha$ is the measured bearing relative to the robot.

The predicted measurement $\mathbf{\hat{z}}_k$ from the robot's predicted state $\mathbf{\hat{x}}_{r,k}$ to the landmark $\mathbf{m}_i$ is:
$$
\mathbf{\hat{z}}_k = \mathbf{h}(\mathbf{\hat{x}}_{r,k}, \mathbf{m}_i) = \begin{bmatrix} \sqrt{\Delta x^2 + \Delta y^2} \\ \operatorname{atan2}(\Delta y, \Delta x) - \hat{\theta}_k \end{bmatrix}
$$
where $\Delta x = m_x - \hat{x}_k$ and $\Delta y = m_y - \hat{y}_k$.

The observation Jacobian $\mathbf{H}_k$, which relates the state change to the measurement change, is computed as the partial derivatives of $\mathbf{h}$ with respect to the state vector. This Jacobian is used to calculate the **Kalman Gain** $\mathbf{K}_k$:
$$
\mathbf{K}_k = \mathbf{\hat{P}}_k \mathbf{H}_k^T (\mathbf{H}_k \mathbf{\hat{P}}_k \mathbf{H}_k^T + \mathbf{R}_k)^{-1}
$$
The state and covariance are then updated:
$$
\mathbf{x}_k = \mathbf{\hat{x}}_k + \mathbf{K}_k (\mathbf{z}_k - \mathbf{\hat{z}}_k)
$$
$$
\mathbf{P}_k = (\mathbf{I} - \mathbf{K}_k \mathbf{H}_k) \mathbf{\hat{P}}_k
$$
The absolute positioning information provided by the RF beacons effectively constrains the scale and eliminates the global drift component accumulated during the VO prediction steps.

% --- Renumbering old sections
\section{Results and Discussions} % Old Section V (now VI)


The hybrid VO–EKF SLAM system was rigorously evaluated in a simulated disaster environment, assessing \textbf{Trajectory Accuracy (ATE/RPE)}, \textbf{Drift Accumulation}, and \textbf{Robustness} under degradation.

\imageplaceholder{image1.png}{Raw VO Trajectory}
\imageplaceholder{image2.png}{Beginning Of Simulation 2D Mapping using EFK}
\imageplaceholder{image3.png}{End Of Simulation 2D Mapping using EFK}
\imageplaceholder{image4.png}{Footage 1: Map Generation Using OpenCV  }
\imageplaceholder{image5.png}{Footage 1: Map Generation With Low Pass Filter Using OpenCV}
\imageplaceholder{image6.png}{Footage 2: Map Generation Using OpenCV}
\imageplaceholder{image7.png}{Footage 2: Map Generation With Low Pass Filter Using OpenCV}
\imageplaceholder{image8.png}{EKF Slam Position Error And Uncertainity Over Time}
\imageplaceholder{image9.png}{Visual Scan Odometry Errors Over Time}

\subsection{Feature Detection and Tracking}
The Shi–Tomasi corner detector yielded an average of $N=350$ stable features per frame. The Lucas–Kanade tracker successfully tracked approximately $90\%$ of the features using a 3-level image pyramid.

\subsection{Essential Matrix Estimation}
RANSAC was used within the Essential Matrix estimation process, filtering out $5\%$ to $15\%$ outlier correspondences, ensuring that the relative motion calculation is based only on the static background structure.

\subsection{Raw VO Trajectory: The Drift Problem}
The trajectory generated solely by integrating the VO motion estimates accumulated a translational error of $\mathbf{8.5\%}$ of the total distance traveled, confirming that the raw VO is not suitable for long-term navigation due to unrecoverable scale factor error.

\subsection{EKF Landmark Corrections}
The EKF correction phase, triggered upon detection of RF beacons, provided vital absolute position updates. This mechanism successfully corrected the scale ambiguity and global orientation, illustrating the 'snap-back' effect where the covariance uncertainty is locally minimized.

\subsection{Hybrid Fused Trajectory and Drift Analysis}
The Hybrid VO–EKF Trajectory achieved an Absolute Trajectory Error (ATE) of less than $0.8\%$ of the total distance traveled, a reduction factor of $\mathbf{10.2}$ compared to the VO-only trajectory. The system utilizes the best features of both methods: VO for precise short-term prediction and EKF-based fusion for global drift correction.

\begin{table}[h]
\caption{Drift Performance Comparison}
\centering
\begin{tabular}{lcc}
\toprule
\textbf{Method} & \textbf{Max Translational Drift} \\
\midrule
VO-only & $8.5\%$ \\
Odometry-only & $15.2\%$ \\
\textbf{Hybrid VO+EKF} & $\mathbf{0.8\%}$ \\
\bottomrule
\end{tabular}
\end{table}

\>
\>
% ---
\section{Robustness and Disaster Constraint Analysis} % Old Section VI (now VII)

\subsection{Occlusion and Low Texture Resilience}
In disaster scenarios, transient occlusion (e.g., debris falling, dust bursts) is common. The EKF framework demonstrates inherent resilience to short-term visual failures. When VO fails to provide an estimate, the EKF relies on the previous state and the propagated $\mathbf{Q}_k$ noise, effectively \textbf{predicting the robot's pose} through the failure period until tracking resumes. This prediction is far superior to relying on uncorrected wheel odometry.

\subsection{Sensor Failure Management}
The multi-modal nature of the EKF allows graceful degradation. If the laser scanner is blinded by smoke (high noise in $\mathbf{R}_{laser}$), the EKF automatically prioritizes the VO and RF beacon measurements. Conversely, if the visual scene is highly textureless, the EKF relies more heavily on the geometric constraints provided by laser scan matching. This dynamic weighting is achieved through the **Kalman Gain** ($\mathbf{K}_k$) calculation.

\subsection{Computational Performance}
The system optimizes computational load by running the high-cost VO ($\approx 30-60$ Hz) only for the prediction step, which is locally precise. The computationally intensive global correction (RF beacons, ICP) runs at a much lower frequency ($\approx 1-5$ Hz) or is event-triggered, enabling real-time operation on resource-constrained robot platforms typical in search-and-rescue.

% ---
\section{Conclusion and Future Work} % Old Section VII (now VIII)

This expanded study provides a full-scale analysis of a hybrid VO–EKF SLAM system for disaster response robotics. VO provides fast, camera-based motion estimation, while the EKF ensures globally consistent localization by fusing VO predictions with auxiliary sensor measurements, crucially including absolute position from RF beacons to correct for scale ambiguity and cumulative trajectory drift.

The system is resilient, cost-efficient, and suitable for field deployment, achieving an ATE of less than $1\%$ of the path length. Future work will focus on three key areas:

\begin{enumerate}
    \item Tighter Visual–Inertial Fusion (VIO): Integrating a low-cost \textbf{Inertial Measurement Unit (IMU)} into the prediction step, potentially using a filter structure like the Multi-State Constraint Kalman Filter (MSCKF), will further enhance short-term prediction, bound the attitude uncertainty, and provide a direct estimate for the unknown scale factor in the VO pipeline.
    \item Loop Closure Integration: Implementing a keyframe-based backend (e.g., pose graph optimization) will provide non-linear, global consistency checks when the robot revisits a previously mapped area. This is essential for long-duration missions where drift accumulation cannot be perfectly corrected by beacons alone.
    \item Real-World Validation: Conducting extensive deployment and validation in accredited disaster training environments (e.g., FEMA sites) to characterize performance against real-world dust, heat, and structural instability.
\end{enumerate}

% ---
\section*{Acknowledgment}
The author thanks the Department of Electrical and Electronics Engineering for providing a challenging environment/deadlines that ultimately facilitated the extended completion time for this project. 

% ---
\begin{thebibliography}{00}

\bibitem{Juneja2014} J. Juneja, ``Scalable Localisation and Mapping of Textured Scenes,'' M.Eng. Thesis, Oxford, 2014.

\bibitem{Juneja2013} J. Juneja, ``Simultaneous Localisation and Mapping (SLAM) in MATLAB,'' 2013.

\bibitem{Triggs2000} B. Triggs et al., ``Bundle Adjustment – A Modern Synthesis,'' Springer, 2000.

\bibitem{Thrun2005} S. Thrun et al., \textit{Probabilistic Robotics}, MIT Press, 2005.

\bibitem{scaramuzza2011} D. Scaramuzza, F. Fraundorfer, ``Visual Odometry,'' IEEE R\&A Magazine, 2011.

\bibitem{bailey2006} T. Bailey, H. Durrant-Whyte, ``SLAM: Part I,'' IEEE R\&A Magazine, 2006.

\bibitem{mur2015} R. Mur-Artal, J. D. Tardos, ``ORB-SLAM: A Versatile and Accurate Monocular SLAM System,'' IEEE TRO, 2015.

\bibitem{engel2014} J. Engel et al., ``LSD-SLAM: Large-Scale Direct Monocular SLAM,'' ECCV, 2014.

\bibitem{engel2017} J. Engel et al., ``Direct Sparse Odometry,'' TPAMI, 2017.

\bibitem{murphy2014} R. Murphy, \textit{Disaster Robotics}, MIT Press, 2014.

\bibitem{tadokoro2009} S. Tadokoro, \textit{Rescue Robotics}, Springer, 2009.

\bibitem{queralta2020} J. Queralta et al., ``Collaborative SLAM in Harsh Environments: A Review,'' IEEE ICMR, 2020.

\bibitem{qin2018} T. Qin et al., ``VINS-Mono: A Robust and Versatile Monocular Visual-Inertial State Estimator,'' IEEE TRO, 2018.

\bibitem{cadena2016} C. Cadena et al., ``Past, Present, and Future of Simultaneous Localization and Mapping: Toward the Robust-Perception Age,'' RA-L, 2016.

\bibitem{forster2014} C. Forster et al., ``SVO: Semi-Direct Visual Odometry for Monocular Cameras,'' ICRA, 2014.

\end{thebibliography}

\end{document}